\section{The exact Perron ground-state transform}
\label{sec:ground}

\begin{theorem}[Dominated Hermitian sum of squares]\label{thm:ground}
Assume \(G_P\) is connected.  For \(\sigma\in\{\pm1\}\) and
\(x_i=r_if_i\),
\begin{equation}\label{eq:ground}
\boxed{\begin{aligned}
x^*(\rho I-\sigma H)x
={}&
\sum_{\{i,j\}\in E(G_P)}r_ir_j
\Bigl[
q_{ij}|f_i-\sigma\phi_{ij}f_j|^2\\
&\hspace{5.0em}
+(p_{ij}-q_{ij})(|f_i|^2+|f_j|^2)
\Bigr].
\end{aligned}}
\end{equation}
On an edge with \(q_{ij}=0\), the first term is interpreted as zero.
\end{theorem}

\begin{proof}
For one undirected edge, the summand expands to
\[
 p_{ij}r_ir_j(|f_i|^2+|f_j|^2)
 -
 2\sigma q_{ij}r_ir_j
 \Real(\overline{f_i}\phi_{ij}f_j).
\]
Summing its diagonal part over edges and using \eqref{eq:perron} gives
\[
 \sum_i r_i|f_i|^2\sum_jp_{ij}r_j
 =
 \rho\sum_i r_i^2|f_i|^2.
\]
The off-diagonal part is exactly \(-\sigma x^*Hx\), because \(H\) is
Hermitian.  This proves \eqref{eq:ground}.
\end{proof}

\begin{corollary}[Exact two-sector variational formula]\label{cor:variation}
Define
\begin{align}
 \cE_\sigma(f)
 =
 \sum_{\{i,j\}}r_ir_j
 \Bigl[
 q_{ij}|f_i-\sigma\phi_{ij}f_j|^2
 +(p_{ij}-q_{ij})(|f_i|^2+|f_j|^2)
 \Bigr].
 \label{eq:energy}
\end{align}
Then
\begin{equation}\label{eq:variation}
 \boxed{\displaystyle
 \delta_\sigma(P,H)
 =
 \min_{f\ne0}
 \frac{\cE_\sigma(f)}
      {\sum_i r_i^2|f_i|^2}.}
\end{equation}
In particular, both sector defects are nonnegative and
\(\|H\|\le\rho(P)\).
\end{corollary}

\begin{proof}
Apply the Rayleigh principle to \(\rho I-\sigma H\) and substitute
\(x_i=r_if_i\).  The substitution is bijective because every
\(r_i>0\).
\end{proof}

\begin{remark}[Two independent sources of energy]
The first term in \eqref{eq:energy} is phase or magnetic frustration.
The second is amplitude-majorant slack.  They are separately
nonnegative; neither can cancel the other.
\end{remark}
